\chapter{Clarifying the ortholog conjecture}

%%%%%%%%%%%%%%%%%%%%%%
%%%%%%%%%%%%%%%%%%%%%%
\section{introduction}
%%%%%%%%%%%%%%%%%%%%%%
%%%%%%%%%%%%%%%%%%%%%%

Othology is a commonly misused term in biology, often referring not to the evolutionary relationship between genes but often being used as a proxy to denote a pair of sequences with high sequence or functional similarity.


most times daughter is not clear (slippage)

neofunctionalization violates conjecture

sub functionalization in some cases does, and does not violate

dosage: unclear

flat curves support cellcular context

age to the last common ancestor is most important for a pair of genes


\section{Clarification of the ortholog conjecture}

The underlying logic behind the ortholog conjecture is  the notion that intra-genomic duplications free one or both copies of the duplicate gene to evolve complementary or novel functions (cite all those papers here).  While the importance of intra-genome duplications is not debated, careful analysis of the potential fates of duplicate genes may not necessarily lead to the same conclusion espoused by the conjecture: that paralogs are more functionally similar than orthologs.

Given the potential fates of duplicate genes outlined by (Matt) I  interpret the ortholog conjecture in the context of stylized gene trees.  I first describe the different trees being considered any why they are important, then go on to consider each potential outcome of a gene duplication.

\section{Stylized gene trees}

\emph{Speciation with no duplication}. This situation is mentioned in order to clarify cases where the ortholog conjecture is untestable.  A  gene family that exhibits no duplication events, by definition will not have any paralogs for use in testing the ortholog conjecture.


\emph{Speciation followed by gene duplication}. I consider two closely related gene tree topologies where speciation precedes duplication and refer to them as SD trees.  This particular topology is especially important because it give rise to one-to-many and many-to-many orthology relationships between organisms, as well as in-paralog relationships within organisms.  \Cref{} shows two such trees, one with only one lineage specific duplication and another with two independent species specific duplication events (referred to as SDD).  These two trees will be especially important in interpreting the ortholog conjecture in the context of subfunctionalization and neofunctionalization.\\[2cm]

\emph{Duplication preceding speciation}. I also consider a gene tree where an intra-genome duplication has occurred before speciation.  I refer to this topology as a DS topology.This topology is important because it gives rise to both within and between species out paralogs.  This is also the topology that is generally considered when discussing very old duplications like those that gave rise to the $\alpha -$ and $\beta -$ hemoglobin chains.  An important point to make is that the ratio of SD to DS trees (SD:DS) among gene families is expected be positively correlated with the age to a pair of organisms' last common ancestor.  A tree showing this topology is shown in \Cref{poo}.  \\

\emph{Duplication and speciation at the same time}.  Finally, I consider trees where the speciation and duplication events occur extremely close to each other or at the same time.  I refer to these topologies as S$=$D topologies.  Although software like Garli \citep{po} and OMA \citep{po} are only capable of labeling a node in a gene tree as both a speciation and duplication node, it could nonetheless be a common topology given the ability of large duplications, or polyploidization events, to lead to symmpatric speciation \citep{po}; or when allopolyploidy occurs.  

\section{The fate of gene duplicates}

Although the ultimate fate of most duplicates is pseudogenization, \citet{po} details several outcome that could befall a conserved gene duplicate.  Here I consider these scenarios and for each stylized topology evaluate the validity of the ortholog conjecture in that context.

\emph{neofuncitonalization}  Neofuncitonalization hypothesizes that, because there are now two copies of a gene, one will be free from evolutionary constraint and capable of developing new functions.  Given a SD topology (\Cref{}) we would expect the ortholog conjecture to hold for the gene copy that retains its original function, but the new copy will be as divergent from its paralog as from its ortholog. In this simplified situation a SD topology would exhibit the ortholog conjecture, but only half of the time.  DS trees would be expected to more strongly support the ortholog conjecture because the duplication occurs before speciation.  A gene would be expected to be more similar to its ortholog than either out-paralog.  

\emph{subfunctionalization}  Subfucntionalization describes a scenario where a duplicate genes that originally carried out two functions will accumulate deleterious mutations that cause each copy of the gene retain complementary ancestral functions.  In an SD tree with only one lineage specific duplication we would have a similar case as with neofunctionalization: paralogs would not exhibit functional similarity, and orthologs would; albiet half the funcitonal similarity as before duplication.  In the second scenario with two independent duplications 




\subsection{Hemoglobin: A case study}

Hemoglobin is one of the most highly cited examples in support of the na\"ive version of the ortholog conjecture \citep{Fitch1970}.  Through its ability to bind iron ions, hemoglobin plays an important role in the transportation of oxygen in a wide range of organisms.  Especially important in utherians is the existence of the fetal $\gamma -$subunit of hemoglobin that exhibits higher oxygen binding affinity and enables the transport of oxygen in the placenta from the.  